\section{Automated Theorem Proving Is Science}
\label{sec:agi}

\subsection{The Critics' Claim}

A recurring objection in discussions of artificial general intelligence runs
as follows. Automated theorem proving---even at superhuman scale, even
across all of mathematics---is not AGI, because formal reasoning is not the
same as doing science. Science requires engaging with the empirical world:
forming hypotheses, running experiments, updating on data. A system that
only manipulates symbols inside a formal system is doing something
categorically different from what scientists do. Theorem proving, however
impressive, is a party trick. It is not intelligence.

This objection is widespread and sounds compelling. We argue it is wrong,
and that it is wrong for a precise reason: it rests on an implicit theory of
science that has never been written down---and when you write it down, it
does not survive scrutiny.

\subsection{The Implicit Theory Being Assumed}

The critics' argument requires a theory of what science is, specifically one
that excludes formal reasoning. The theory being assumed, roughly, is this:
science is the activity of interacting with the physical world through
observation and experiment, and formal reasoning is merely a tool that
science uses, not science itself.

This theory has intuitive appeal but no formal content. It cannot tell you:
\begin{itemize}
  \item When two competing theories are equally consistent with the data,
        which is better?
  \item How much complexity is a theory allowed before it is no longer
        parsimonious?
  \item When should a paradigm be abandoned in favor of a competitor?
  \item What counts as a legitimate auxiliary hypothesis versus an ad hoc
        patch?
\end{itemize}
These are the central questions of the philosophy of science, and the
informal ``interaction with the physical world'' theory has no answers to
them. It is not a theory of science. It is a description of one setting in
which science sometimes happens.

\subsection{The Formal Theory}

The Bayesian/Kolmogorov framework provides an actual theory of science,
with formal answers to all of the above questions. Restating it precisely:

\begin{definition}[Scientific activity]
  A process is engaged in scientific activity if it:
  \begin{enumerate}
    \item Maintains a space of computable hypotheses $\mathcal{H}$;
    \item Assigns prior probabilities $P(H) \propto 2^{-\K(H)}$ to
          hypotheses, preferring shorter descriptions;
    \item Evaluates hypotheses against observations via a likelihood
          function $P(D \mid H)$;
    \item Updates beliefs by Bayes' rule: $P(H \mid D) \propto
          P(D \mid H) \cdot P(H)$; and
    \item Selects hypotheses that minimize $L(H) + L(D \mid H)$.
  \end{enumerate}
\end{definition}

Notice what is \emph{not} in this definition: laboratories, physical
instruments, empirical data in any particular form, or indeed any
requirement that $D$ come from the external world rather than from a
formal system. The definition is about the \emph{structure} of the
inference, not the \emph{substrate} in which it occurs.

\subsection{The Verification Oracle Is Interchangeable}

The key insight is that the search/verify asymmetry
(Section~\ref{sec:verification}) holds regardless of what the verifier is.

In empirical science, the verifier is experiment: you propose a hypothesis,
derive a prediction, run the experiment, and check whether the prediction
holds. Finding a good hypothesis is hard; running the experiment is cheap
relative to the search.

In formal mathematics, the verifier is the proof checker: you propose a
theorem and a proof, and the type-checker verifies it in time linear in
proof length. Finding a good proof is hard (it may require deep insight or
exhaustive search); checking it is cheap.

Both are instances of the same abstract structure:
\[
  \text{Search}(H \in \mathcal{H}) \xrightarrow{\text{expensive}}
  \hat{H} \xrightarrow{\text{cheap verification}} \text{accept or reject}
\]
The oracle that implements ``cheap verification'' differs in its physical
instantiation---a particle accelerator versus a type-checker---but the
computational structure is identical. A system that can search hypothesis
space and verify proposals is doing science. The nature of the oracle is
irrelevant to this classification.

\subsection{What Formal Mathematics Looks Like Under This Theory}

Under the Kolmogorov framework, a mathematical theorem is a hypothesis about
the structure of a formal system. A proof is the verification. The
``observations'' $D$ are the axioms and previously established theorems. The
``experiment'' is type-checking.

This is not a strained analogy. Consider what a mathematician does:
\begin{itemize}
  \item She maintains a space of possible theorems (hypotheses).
  \item She assigns higher prior probability to theorems that are short,
        elegant, and follow from a few strong principles---precisely
        $P(H) \propto 2^{-\K(H)}$.
  \item She evaluates a candidate theorem against the existing body of
        mathematics (the ``data'').
  \item She updates: a theorem that follows from deep principles and
        connects many areas is more valuable than one that requires a
        long, ad hoc proof.
  \item She publishes results that compress the existing body of knowledge
        into shorter descriptions---unifying previously separate results,
        revealing that two seemingly different theorems are instances of
        one.
\end{itemize}
This is the MDL objective, applied to a formal domain. The mathematician
is doing exactly what the physicist does. The substrate differs; the
objective function is the same.

\subsection{The Objection Answered}

The critics are right that theorem proving, as typically practiced, is not
the same as doing science. But this is because automated theorem provers as
currently deployed do not have the right objective function. They are
optimized to prove specified theorems, not to search for theorems worth
proving. They are given $H$ and asked to verify it, rather than searching
$\mathcal{H}$ for the $H$ that minimizes description length.

This is an engineering gap, not a conceptual impossibility. A theorem-proving
system equipped with:
\begin{enumerate}
  \item A prior over theorems weighted by description length,
  \item A likelihood that measures how much a theorem compresses and
        unifies the existing mathematical corpus,
  \item A search procedure over the space of provable statements, and
  \item A cheap verifier (the type-checker),
\end{enumerate}
would be doing science in the full formal sense. It would propose conjectures
that compress the mathematical corpus, verify them, keep the ones that
advance the MDL objective, and discard the rest. This is Solomonoff
induction instantiated in a formal domain.

Such a system would not be doing something categorically different from what
Einstein did when he proposed general relativity, or what Ramanujan did when
he produced his notebooks. It would be doing the same thing, faster and
without biological constraints.

\subsection{AGI and the Objective Function}

The question of whether a system constitutes artificial general intelligence
is often posed in terms of capabilities: can it reason, plan, learn, adapt?
We propose a different framing, grounded in the framework of this paper.

\begin{proposition}[AGI as correct objective function]
  A system exhibits artificial general intelligence if and only if it
  implements---exactly or approximately---Solomonoff induction: maintaining
  a prior over computable hypotheses weighted by description length and
  updating by Bayes' rule on available observations.
\end{proposition}

This definition has several virtues. It is formal and precise. It is
substrate-independent: the system can operate on physical observations,
mathematical structures, or any other computable domain. It is graded:
systems can approximate Solomonoff induction more or less closely, with
better approximations constituting more general intelligence. And it
connects AGI directly to the objective function for science: a system that
is doing science in the formal sense is, by this definition, exhibiting
general intelligence.

Under this definition, the question ``is theorem proving AGI?'' has a
precise answer: it depends on whether the theorem-proving system has the
right objective function. A system that is given theorems to prove and
proves them is not AGI---it is a very capable special-purpose tool. A
system that searches for theorems worth proving, verifies them, and updates
its prior based on what it finds---minimizing description length over the
space of mathematical knowledge---is AGI, regardless of whether it ever
runs a physical experiment.

\subsection{The Broader Implication}

The critics who dismiss theorem proving as ``not really AGI'' are, without
knowing it, making an implicit claim about the objective function for
science. They are claiming that the objective function requires physical
observation as input. We have argued that this claim is false: the objective
function is MDL, the input can be any observations from any verifiable
domain, and the structure of inference is the same regardless.

The tools to build such a system exist. The objective function is written
down---in this paper and in the century of work it synthesizes. The
verification infrastructure exists: Lean~4 type-checks proofs in linear
time; benchmark evaluation checks model performance in seconds; experimental
platforms check empirical predictions in hours or days. The search
infrastructure is being built: AutoResearch~\cite{karpathy2026},
agent-based research systems, and large language models used as hypothesis
generators are all approximating the search step.

What has been missing is the theoretical unification: the recognition that
all of these are the same process, that they all approximate the same
objective function, and that a system implementing that objective function
in any verifiable domain is doing science and exhibiting general
intelligence.

That recognition is the contribution of this paper.